\chapter{Review of Literature and Plagiarism Issues}

\section{Historical Background of Malaria Control}

Malaria has plagued humanity for millennia, with descriptions of periodic fevers dating back to ancient Greek and Roman civilizations. The discovery that mosquitoes transmit malaria by \citet{ross1897} and \citet{grassi1900} revolutionized public health approaches to disease control. Early control efforts focused on vector control through environmental management, including draining swamps and spraying insecticides.

The advent of dichlorodiphenyltrichloroethane (DDT) during World War II marked a turning point in malaria control. The Global Malaria Eradication Program (1955-1969), led by the World Health Organization, relied heavily on indoor residual spraying with DDT. While the program eliminated malaria from Europe and North America, it failed in sub-Saharan Africa due to logistical challenges, insecticide resistance, and weak health infrastructure \citep{who2023malaria}.

The resurgence of malaria in the 1970s and 1980s prompted a shift toward integrated vector management. Insecticide-treated nets (ITNs) emerged as a promising intervention in the 1990s, with early trials demonstrating substantial reductions in child mortality \citep{lengeler2004insecticide}.

\section{Evidence for ITN Effectiveness}

\subsection{Randomized Controlled Trials}

The gold standard evidence for ITN effectiveness comes from randomized controlled trials. A Cochrane review by \citet{lengeler2004insecticide} synthesized evidence from 22 trials and found that ITNs reduce child mortality by 17\% (RR 0.83, 95\% CI 0.77-0.89) and reduce episodes of malaria by 50\%. An updated Cochrane review by \citet{pryce2018insecticide} confirmed these findings, reporting a 17\% reduction in \textit{P. falciparum} prevalence (RR 0.83, 95\% CI 0.71-0.98).

\subsection{Observational Studies}

While RCTs provide strong internal validity, their generalizability to real-world settings is limited. Observational studies using large-scale survey data have complemented RCT evidence. \citet{lim2011net} conducted a multicountry analysis of Demographic and Health Survey data from 22 African countries, using matching methods to estimate ITN effectiveness. They found a pooled relative reduction in parasitemia prevalence of 20\% (95\% CI 3-35\%) associated with household ITN ownership.

\citet{bhatt2015effect} estimated that ITNs contributed to a 20-30\% reduction in malaria prevalence in Africa between 2000 and 2015, accounting for approximately 68\% of the total reduction achieved during that period.

\subsection{Evidence from Kenya}

Several studies have examined malaria epidemiology in Kenya. \citet{noor2014risk} mapped malaria risk across Kenya, identifying high-transmission zones in the western region near Lake Victoria and coastal areas. \citet{ayele2024assessment} used multilevel logistic regression to analyze malaria transmission in Kenya, demonstrating the importance of accounting for clustering at the enumeration area level.

The Kenya Malaria Indicator Surveys \citep{kmis2015report, kmis2020report} provide nationally representative data on malaria prevalence and ITN use. These reports show that malaria prevalence among children under five declined from 8.2\% in 2010 to 4.8\% in 2020, while ITN use increased from 40\% to 54\% over the same period.

\section{Causal Inference Methods in Epidemiology}

\subsection{Traditional Methods}

Traditional causal inference methods in epidemiology include logistic regression, which estimates associations between exposures and outcomes while adjusting for confounders. However, logistic regression assumes linearity in the log-odds and may produce biased estimates when confounding is severe \citep{austin2011introduction}.

Generalized Linear Mixed Models (GLMM) extend traditional regression by incorporating random effects to account for clustered data structures. As \citet{neuhaus1991comparison} demonstrate, when intra-cluster correlation is high, cluster-specific (conditional) effects differ substantially from population-average (marginal) effects.

\subsection{Propensity Score Methods}

The propensity score, defined as the probability of treatment given observed covariates, was introduced by \citet{rosenbaum1983central}. The key insight is that conditioning on the propensity score balances observed covariates between treated and untreated groups, mimicking a randomized experiment.

\citet{austin2011introduction} provides a comprehensive review of propensity score methods, including matching, inverse probability weighting, and stratification. \citet{austin2010optimal} established optimal caliper widths for propensity score matching.

Doubly robust estimation \citep{kang2007average} combines propensity score weighting with outcome regression, remaining consistent if either model is correctly specified.

\subsection{Machine Learning for Causal Inference}

Double Machine Learning \citep{chernozhukov2018double} combines machine learning algorithms with causal inference techniques. The method uses Neyman-orthogonal scores and cross-fitting to provide valid inference even when nuisance functions are estimated using flexible machine learning methods.

\citet{nie2021quasi} developed the R-learner, which estimates heterogeneous treatment effects using orthogonalization. \citet{wu2025why} recently identified limitations in DML variance estimation, showing that influence-function-based variance estimators are not doubly robust.

\subsection{The Positivity Assumption}

\citet{westreich2010invited} emphasize the importance of the positivity assumption in causal inference, which requires that every combination of covariates has a non-zero probability of both treatment and non-treatment. This assumption must be empirically verified.

\section{Gaps in the Literature}

Despite extensive research, several gaps remain. First, most observational studies of ITNs report associations rather than causal effects. Second, existing analyses often ignore the hierarchical structure of survey data, despite evidence of substantial clustering (ICC = 41-47\% in Kenya). Third, few studies distinguish between cluster-specific (conditional) and population-average (marginal) effects, which have different interpretations for policy.

This study addresses these gaps by: (1) applying rigorous causal methods, (2) explicitly accounting for clustering using GLMM for propensity score estimation, and (3) estimating both conditional and marginal effects.

\section{Plagiarism and Ethical Considerations}

\subsection{Definition of Plagiarism}

Plagiarism is the act of presenting someone else's work, ideas, or words as one's own without proper attribution. It includes: (1) direct copying without quotation marks and citation, (2) paraphrasing without citation, (3) using images or figures without permission and attribution, and (4) self-plagiarism (reusing one's own previously published work without citation).

\subsection{How This Study Avoids Plagiarism}

In accordance with AIMS academic integrity policies, the following measures were taken to avoid plagiarism:

\begin{itemize}
	\item \textbf{Proper citation}: All sources of information, data, and ideas are cited using the author-date citation style. Each citation in the text corresponds to an entry in the reference list.
	
	\item \textbf{Quotation marks}: Direct quotations of fewer than 40 words are enclosed in quotation marks with page numbers provided. Longer quotations are indented as block quotes.
	
	\item \textbf{Paraphrasing}: All paraphrased content has been rewritten in the author's own words while still citing the original source.
	
	\item \textbf{Original figures and tables}: All figures and tables were generated by the author using R code. Where data from external sources are presented, the source is cited in the caption.
	
	\item \textbf{AI usage disclosure}: The use of generative AI (DeepSeek) for code assistance and grammar checking is disclosed in the Appendix, with supervisor approval documented.
\end{itemize}

\subsection{Declaration of Originality}

The author declares that this work is original and has not been submitted for any other degree or qualification. Any work done by others or by the author previously has been acknowledged and referenced accordingly. The author takes full responsibility for the accuracy, integrity, and originality of the work presented herein.